Assumption 4 (Survival Factor Model). For all \(p\in[P]\), \(n\in[N]\) and \(t\in[0,\widetilde{\tau}]\), the potential survival function under control \(S^{(0)}_{p,n}(t)=\langle u_p(t),v_n\rangle\), where \(u_p:[0,\widetilde{\tau}]\to\mathbb R^r\) are (latent) period factors and \(v_n\in\mathbb R^r\) are (latent) unit factors for some \(r\ge1\). Assumption 6 (Linear Span Inclusion). For each treated unit \(n\in\mathcal I^{(1)}\), there exists weights \(w_{n,m}\in\mathbb R^{|\mathcal I^{(0)}|}\) such that \(v_n=\sum_{m\in\mathcal I^{(0)}}w_{n,m}\cdot v_m\). Theorem 1. Given the assumptions above, we have \(S^{(0)}_{1,n}(\cdot)=\sum_{m\in\mathcal I^{(0)}}w_{n,m}S^0_{1,m}(\cdot)\). Assumption 7 (Well-conditioned Spectrum). The expectation of the control group pre-treatment matrix satisfies: \(\operatorname{cond}(\mathcal S_{0,\mathcal I^{(0)}})\le\xi'\), and \(\|\mathcal S_{0,\mathcal I^{(0)}}\|_F^2\ge\xi''N_0T_0\) for some constants \(\xi',\xi''>0\), where \(\|\cdot\|_F\) is the Frobenius norm. Assumption 8 (Latent time-period span). We assume the post-period specific factors lie in the span of pre-period specific factors, i.e., \(\forall t\in\mathcal T,u_1(t)\in\operatorname{span}(\{u_0(s):s\in\mathcal T\})\) conditional on \(\{\mathcal{LF},\mathcal D\}\).